\begin{document}

\section{提案手法}

\subsection{パイプライン}

図\ref{fig:pipeline}に，本研究で用いるCLOUSEAU型のエージェント型ログ探索パイプラインを示す．入力は，CBC alert rows，またはhost/process/time windowである．CLOUSEAU本来のPOI起点調査\cite{clouseau}にならい，Chief agentは初期手がかりから調査方針を生成する．Investigation agentは，確認すべき事実をQAAgent / SQL expertへの質問に変換する．QAAgent / SQL expertは，SQLite化されたエンドポイントログに問い合わせ，timestamp，source stream，PID/PPID，process tree，command line，registry path，source row idなどを返す．Final synthesisは観測結果を統合し，主行動列，補助証跡，limitationsを出力する．

\begin{figure*}[t]
\centering
\fbox{\begin{minipage}{0.92\linewidth}
\small
\centering
SOC clue / CBC alert / process-time clue\\
$\downarrow$\\
Chief agent: 初期手がかりから調査方針を生成\\
$\downarrow$\\
Investigation agent: 確認すべき事実を質問へ変換\\
$\downarrow$\\
QAAgent / SQL expert: SQLite endpoint logsから観測行を取得\\
$\downarrow$\\
Final synthesis: 主行動列，証跡，近傍行動，limitationsを統合\\
$\downarrow$\\
component rubricで採点
\end{minipage}}
\caption{SOC調査起点から証跡付き行動列を復元するパイプライン}
\label{fig:pipeline}
\end{figure*}

\subsection{出力と評価単位}

本研究では，最終出力を行動ステップ単位で評価する．各gold stepは，主体，操作，対象からなるaction componentと，そのstepをログから支えるcritical evidenceを持つ．さらに，隣接するgold stepの順序関係をsequence orderとして評価する．モデル出力がgold内容を実質的に含めばhitとし，表現の違いは許容する．一方，alert nameやalert reasonを行動そのものとして出した場合，gold chainに属さない近傍行動を主系列として含めた場合，証跡で支えられない推定を含めた場合は，candidate claim precisionの分母またはoverclaimとして扱う．

\end{document}
